
\documentclass[10pt]{article}

\usepackage[letterpaper,top=1in,left=1in,right=1in]{geometry}
\usepackage[dvipsnames,svgnames]{xcolor}
\usepackage[shortlabels]{enumitem}
\usepackage[framemethod=TikZ]{mdframed}
\usepackage{amsmath,amssymb,amsthm}
\usepackage[colorlinks]{hyperref}
\usepackage{multicol}
\usepackage{tabularx}
\usepackage{bbm}
\usepackage{tikz-cd}
\usepackage{mathrsfs}

\hypersetup{
    colorlinks=true,
    linkcolor=blue,
    filecolor=magenta,
    urlcolor=magenta,
    pdftitle={{CA comps problems}}
}

\usepackage{mathtools}
\usepackage{thmtools}
\usepackage[textsize=footnotesize]{todonotes}
\usepackage{titling}
\usepackage{cancel}

\mdfdefinestyle{mdblackbox}{linecolor=black,backgroundcolor=RedViolet!5!gray!5,
linewidth=3pt,rightline=false,leftline=true,topline=false,bottomline=false,}
\declaretheoremstyle[mdframed={style=mdblackbox}]{thmblackbox}

\declaretheoremstyle[spaceabove=6pt,spacebelow=6pt]{boldhead}

\declaretheorem[style=boldhead,name=Problem]{problem}
\declaretheorem[style=boldhead,name=Claim,numberwithin=problem]{claim}
\declaretheorem[style=thmblackbox,name=Remark,numberwithin=problem]{remark}

\providecommand{\ol}{\overline}
\providecommand{\tensor}{\otimes}
\renewcommand{\hom}{\operatorname{Hom}}
\providecommand{\tor}{\operatorname{Tor}}
\providecommand{\Gal}{\operatorname{Gal}}
\providecommand{\ceil}[1]{\left\lceil #1 \right\rceil}
\providecommand{\floor}[1]{\left\lfloor #1 \right\rfloor}
\newcommand{\dd}{\,\mathrm{d}}
\providecommand{\ul}{\underline}
\providecommand{\eps}{\varepsilon}
\providecommand{\half}{\frac{1}{2}}
\providecommand{\CC}{\mathbb C}
\providecommand{\FF}{\mathbb F}
\providecommand{\II}{\mathbb I}
\providecommand{\NN}{\mathbb N}
\providecommand{\QQ}{\mathbb Q}
\providecommand{\RR}{\mathbb R}
\providecommand{\ZZ}{\mathbb Z}
\providecommand{\ind}{\mathbbm 1}
\providecommand{\dg}{^\circ}
\providecommand{\ii}{\item}
\providecommand{\setdiff}{\smallsetminus}
\providecommand{\alert}[1]{{\sffamily\textbf{\textcolor{blue}{#1}}}}
\providecommand{\inv}{^{-1}}
\providecommand{\mb}[1]{\mathbf{#1}}
\newcommand{\what}{\widehat}

\newenvironment{soln}{\begin{proof}[Solution]}{\end{proof}}

\providecommand{\pow}{\operatorname{Pow}}
\providecommand{\vocab}[1]{{\textbf{\textcolor{ForestGreen}{#1}}}}
\providecommand{\lcm}{\operatorname{lcm}}
\providecommand{\abs}[1]{\left\lvert #1\right\rvert}
\providecommand{\smabs}[1]{\lvert #1\rvert}
\providecommand{\innerprod}[1]{\langle\!\langle #1\rangle\!\rangle}
\providecommand{\norm}[1]{\left\| #1\right\|}
\providecommand{\smnorm}[1]{\| #1\|}
\providecommand{\gr}{\operatorname{gr}}

\begin{document}

\title{SBU Comps Complex Analysis Problems}
\author{Aditya Pahuja}
\date{\today}
\maketitle

\begin{problem}
[2026 J-I-3]
Prove that for every integer $n\geq 0$ and every entire function $f$ on the complex plane, $f$ is a degree-$n$ polynomial if and only if
$$
\lim_{z\to\infty}\frac{zf'(z)}{f(z)}=n.
$$
\end{problem}

\begin{soln}
    If $f$ is a degree $n$ polynomial
    then $zf'(z)$ is also a degree $n$ polynomial
    whose leading coefficient is $n$ times
    the leading coefficient of $f$ (zero in the case where $n=0$),
    so obviously
    \begin{equation*}
	\lim_{z\to\infty}\frac{zf'(z)}{f(z)} = n.
    \end{equation*}

    Suppose now that the limit condition holds.
    The nonzero poles of $\frac{zf'(z)}{f(z)}$ are
    exactly the nonzero zeroes of $f(z)$,
    so in order for the given limit to be finite,
    the set of zeroes must be bounded
    (since the magnitude of a meromorphic function
    near a pole is unbounded).
    Thus $f\inv(0)$ is closed in $\CC$ and bounded,
    hence compact, and it is discrete since $f$ is obviously
    not the zero function; this means $f\inv(0)$ is finite.

    Suppose that the zeroes of $f$,
    with multiplicity, are $z_1$, $z_2$, \dots, $z_m$.
    Then, $f(z) = (z - z_1)\cdots(z - z_m)g(z)$
    for a nowhere vanishing entire function $g$,
    which means that
    \begin{equation*}
	\frac{zf'(z)}{f(z)}
	= \frac{z}{z - z_1} + \frac{z}{z - z_2}
	+ \cdots + \frac{z}{z - z_m}
	+ \frac{zg'(z)}{g(z)}.
    \end{equation*}
    Since $g$ is nowhere zero we can write
    $g(z) = e^{h(z)}$ for some entire function $h$.
    Taking the limit gives
    \begin{equation*}
	n = \lim_{z\to\infty} \frac{zf'(z)}{f(z)}
	= m + \lim_{z\to\infty} \frac{zg'(z)}{g(z)}
	= m + \lim_{z\to\infty} zh'(z),
    \end{equation*}
    so $zh'(z)$ is bounded and entire, hence
    $zh'(z) = 0\cdot h'(0) = 0$ for all $z\in\CC$.
    This implies the limit in the right-hand side is zero,
    i.e. $m = n$, and $h$ and $g$ are constant functions,
    say $g(z) = c$ for all $z$ and nonzero $c$
    (if $c$ were zero then the given limit would not make sense),
    so we have shown that $f(z) = c(z - z_1)\cdots(z - z_n)$.
\end{soln}



\begin{problem}
[2026 J-II-4]
Let $U$ be an open subset of $\mathbb{C}$. Let $(f_j(z))$ be a collection of holomorphic functions on $U$ that converges uniformly on every compact subset. Carefully prove that the pointwise limit is holomorphic.
\end{problem}

\begin{soln}
    Is this not just showing that $f$ is holomorphic
    around any $z_0\in U$ by Morera's theorem?
\end{soln}



\begin{problem}
[2025 A-I-3]
Among all monic polynomials of degree $n$, show that $z^n$ minimizes
$$
\sup_{|z|<1}|p(z)|.
$$
\end{problem}

\begin{soln}
    The maximum magnitude of $p(z)$ on the interior of the unit disk
    is equal to the maximum along the disk's boundary,
    by continuity and the maximum modulus principle.
    Write $p(z) = a_nz^n + a_{n-1}z^{n-1} + \cdots + a_1z + a_0$
    for $a_0,\dots,a_n\in\CC$ and $a_n = 1$.
    Then $\abs{p(z)}^2$ can be written for $z$ on the unit circle as
    \begin{equation*}
	p(z)\cdot\ol{p(z)}
	= \sum_{j=0}^n\sum_{k=0}^n a_j\ol{a_k}\cdot z^{j-k},
    \end{equation*}
    so the average of $\abs{p(z)}^2$ on the unit circle is
    \begin{equation*}
	\frac 1{2\pi}\int_0^{2\pi}
	p(e^{i\theta})\cdot\ol{p(e^{i\theta})}\dd\theta
	= \frac 1{2\pi}\sum_{0\le j,k\le n}a_j\ol{a_k}\int_0^{2\pi}
	e^{i\theta(j-k)}\dd\theta.
    \end{equation*}
    The integral of $e^{i\theta(j-k)}$ is zero when $j\ne k$
    because the real and imaginary parts are
    $\cos((j-k)\theta)$ and $\sin((j-k)\theta)$,
    both of which integrate to zero over their period when $j-k\ne 0$,
    and $2\pi$ when $j = k$, since in that case
    the integrand is just $1$. That is, the average of $\abs{p(z)}^2$
    on the unit circle is
    \begin{equation*}
	\frac 1{2\pi}\int_0^{2\pi}
	p(e^{i\theta})\cdot\ol{p(e^{i\theta})}\dd\theta
	= \frac 1{2\pi}\sum_{j=0}^n a_j\ol a_j
	= 1 + \abs{a_{n-1}}^2 + \cdots + \abs{a_1}^2 + \abs{a_0}^2
	\ge 1,
    \end{equation*}
    so for any monic polynomial,
    $\sup_{\abs{z} < 1}\abs{p(z)}$ is at least $1$,
    with equality (corresponding to the case where
    almost every point on the boundary has magnitude $1$)
    if and only if each of $a_0$, \dots, $a_{n-1}$ are $0$,
    i.e. $p(z) = z^n$.
\end{soln}



\begin{problem}
[2025 A-II-3]
Let $f:\mathbb{D}\to U$ be a conformal map from the unit disk $\mathbb{D}=\{z:|z|<1\}\subset\mathbb{C}$ to a domain $U\subset\mathbb{C}$ of bounded area. For $\theta\in[0,2\pi)$, let $R_\theta=\{re^{i\theta}:1/2\leq r<1\}$. Show that
$$
\int_0^{2\pi}\text{length}(f(R_\theta))\,d\theta<\infty.
$$ Show that for almost any angle $\theta$, there exists the radial limit
$$
\lim_{r\to 1}f(re^{i\theta}).
$$
\end{problem}

\begin{soln}

\end{soln}



\begin{problem}
[2025 J-I-4]
Find all entire functions $f$ and $g$ satisfying the equation
$$
e^f+e^g=4.
$$
\end{problem}

\begin{soln}
    Since $e^f$ is never zero, $e^g$ can never be equal to $4$,
    and vice versa. By Picard's theorem this means $e^f$ and $e^g$,
    hence $f$ and $g$, are both constant
    (they are entire and miss the values $0$ and $4$).
    Thus the only functions satisfying this equation are
    $f\equiv a$, $g\equiv b$ for any pair of consants $a,b\in\CC$
    satisfying $e^a + e^b = 0$.
\end{soln}



\begin{problem}
[2025 J-II-1]
Let $f\colon\mathbb{D}\to\mathbb{D}$ be a nonconstant holomorphic map such that for every compact subset $K\subset\mathbb{D}$, the set $f^{-1}(K)$ is compact. Show that there is a number $N\in\mathbb{N}$ such that for every $a\in\mathbb{D}$, the set $f^{-1}(a)$ contains at most $N$ distinct points in $\mathbb{D}$.
\end{problem}

\begin{soln}
    Since $f$ is proper, $f\inv(a)$ is finite for each $a\in\mathbb D$,
    say with $n_a$ elements. Then $f(z) - f(a)$ has
    $n$ zeroes (with multiplicity).
\end{soln}



\begin{problem}
[2024 J-I-3]
Suppose $f$ is continuous on the plane and holomorphic on $\mathbb{C} \backslash \mathbb{R}$. Prove that $f$ is holomorphic on the whole plane.
\end{problem}

\begin{soln}
    Morera
\end{soln}



\begin{problem}
[2023 A-II-6]
Find the number of solutions (counting multiplicity) to $z^8-5 z^6+2 z^3-z-1=0$ that lie inside the unit disk.
\end{problem}

\begin{soln}
    If $z$ is a solution to the equation with $\abs z = 1$,
    \begin{equation*}
	\smabs{5z^6} = \smabs{z^8 + 2z^3 - z - 1}
	\le \smabs{z^8} + \smabs{2z^3} + \smabs{z} + \smabs{1}
	= 5 = \smabs{5z^6},
    \end{equation*}
    so equality holds in the triangle inequality.
    This happens if and only if $1$, $z$, $2z^3$, and $z^8$
    are real multiples of one another,
    i.e. $z$ is one of $\pm 1$, but $z = 1$
    and $z = -1$ are easily verified to not be solutions.
    Thus,
    \begin{equation*}
	\smabs{z^8 + 2z^3 - z - 1} < \smabs{5z^6}
    \end{equation*}
    so Rouch\'e's theorem implies that
    the given equation has the same number of zeroes in
    the unit disk as $5z^6$, namely $\boxed 6$.
\end{soln}



\begin{problem}
[2023 A-I-6]
Let $\mathbb{D} \subset \mathbb{C}$ be the open unit disk. Is there a holomorphic function $f$ with $f(\mathbb{D})=\mathbb{D}$, $f(0)=f^{\prime}(0)=2 / 3$? If so, give a formula. If not, prove that it cannot exist.
\end{problem}

\begin{soln}

\end{soln}



\begin{problem}
[2023 J-I-2]
Set $\mathbb{D}=\{z \in \mathbb{C} :  \smabs z<1\}$ and let $f: \mathbb{D} \to \{w \in \mathbb{C} : e^{-\pi / 2}<|w|< e^{\pi / 2}\}$ be a holomorphic map satisfying $f(0)=1$. Show that $\left|f^{\prime}(0)\right| \leq 2$.
\end{problem}

\begin{soln}

\end{soln}



\begin{problem}
[2022 A-I-5]
If $\Omega \subset \mathbb{C}$ is simply connected and $f: \Omega \rightarrow \mathbb{C} \backslash\{0\}$ is a holomorphic function, is there a holomorphic function $g: \Omega \rightarrow \mathbb{C}$ such that $f=\exp (g)$ on $\Omega$? Prove or give a counterexample.
\end{problem}

\begin{soln}

\end{soln}



\begin{problem}
[2022 A-II-3]
Let $\Omega$ be the half-disk $\{z=x+i y: x>0,|z|<1\}$. Let $u$ be the bounded harmonic function on $\Omega$ that has boundary value 1 on the vertical open segment $(-i, i)$ and has boundary value 0 on the open circular arc $\{|z|=1, x>0\}$. What is the value of $u(1 / 2) ?$
\end{problem}

\begin{soln}

\end{soln}



\begin{problem}
[2022 J-I-1]
Suppose $f: \mathbb{D} \longrightarrow \mathbb{C}$ is a holomorphic function on the unit disk so that $$f\left(\frac{1}{n}\right)=\frac{5}{n^2} \quad \forall n \in \mathbb{Z}, n \geq 2.$$ Find $f^{\prime \prime}(0)$.
\end{problem}

\begin{soln}
    By continuity we have that $f(0) = 0$,
    so the zero set of $f(z) - 5z^2$ is not discrete,
    meaning that $f(z) - 5z^2$ is the zero function.
    Thus $f''(z) = 10$ for all $z$, especially $z = 0$.
\end{soln}



\begin{problem}
[2022 J-II-4]
Let $U \subset \mathbb{C}$ be an open subset. Suppose $f_i: U \longrightarrow \mathbb{C}$ is a sequence of holomorphic functions converging uniformly on compact subsets to a function $f: U \longrightarrow \mathbb{C}$. Show that $f$ is also holomorphic. Justify each step clearly.
\end{problem}

\begin{soln}
    This problem is the same as 2026 J-II-4.
\end{soln}



\begin{problem}
[2021 A-I-2]
Let $P(z)=z^n+a_1 z^{n-1}+\cdots+a_n$. Prove that either there is $|z|=1$ such that $|P(z)|>1$ or $P(z)=z^n$.
\end{problem}

\begin{soln}
    This problem is the same as 2025 A-I-3.
\end{soln}



\begin{problem}
[2021 J-I-6]
By Riemann's mapping theorem there is a conformal map $f:(-1,1)^2 \rightarrow$ $D$ from the open square into the open disc $D=\{z:|z|<1\}$ in the complex plane which is uniquely determined by the conditions $f(0)=0$ and $f^{\prime}(0)>1$. Show that the radius of convergence of the Taylor series for $f$ at the origin is less than or equal to 2.
\end{problem}

\begin{soln}

\end{soln}



\begin{problem}
[2020 A-I-5]
Let $\Omega$ be the half-disk $\{z=x+i y : x>0, \smabs z<1\}.$ Let $u$ be the bounded harmonic function on $\Omega$ that has boundary values 1 on the vertical open segment $(-i, i)$ and has boundary values 0 on the open circular arc $\{|z|=1, x>0\}.$ What is the value of $u(1 / 2) ?$
\end{problem}

\begin{soln}

\end{soln}



\begin{problem}
[2020 A-II-1]
How many roots (counted with multiplicity) does the function
$$
g(z)=6 z^3+e^z+1
$$
have in the unit disk $|z|<1 ?$
\end{problem}

\begin{soln}
    Since
    \begin{equation*}
        \smabs{e^z + 1} \le e + 1 < 6 = \smabs{6z^3}
    \end{equation*}
    Rouch\'e's theorem says that $g(z)$ and $6z^3$
    have the same number of solutions on the unit disk,
    namely $\boxed 3$.
\end{soln}



\begin{problem}
[2020 J-II-3]
Let $\mathcal{H}=\{x+i y \mid y >0\} \subset \mathbb{C}$ be the upper half-plane, and let $\mathcal{U}=\mathcal{H}-\gamma$, where $\gamma$ is the line segment $\{3+i y \mid 0< y \leq 1\}$. Construct an explicit holomorphic diffeomorphism $\Phi: \mathcal{U} \rightarrow \mathcal{D}$, where $\mathcal{D}$ is the open unit disk. You may leave your answer expressed as a composition of other maps.
\end{problem}

\begin{soln}

\end{soln}



\begin{problem}
[2019 J-I-2]
Evaluate the integral
$$
\int_0^{\infty} \frac{x^{a-1}}{x+1} d x
$$
where $0< a<1$.
\end{problem}

\begin{soln}

\end{soln}



\begin{problem}
[2019 J-II-3]
Denote by $D_\rho$ the disk of radius $\rho$ with center at 0 on the complex plane $\mathbb{C}$. For each $R>0$ and each $p \in D_R$ let
$S_R(p) = \sup\{r > 0 : \exists f\colon D_r\to D_R\text{ holomorphic, }
f(0) = p, \smabs{f'(0)} = 1\}$.
Find $S_R(p)$.
\end{problem}

\begin{soln}

\end{soln}



\begin{problem}
[2019 A-II-1]
Suppose $f: \mathbb{C} \rightarrow \mathbb{C}$ is continuous on the closed disk $\{z:|z| \leq 1\} \subset \mathbb{C}$ and holomorphic on its interior. If $f$ restricted to the boundary of this disk only takes values on the curve $x y=1$ (where a point $z \in \mathbb{C}$ is written as $z=x+i y,$ with $x, y \in \mathbb{R}$ ), show that $f$ is constant on the whole disk.
\end{problem}

\begin{soln}

\end{soln}



\begin{problem}
[2018 J-I-2]
Use the method of residues to compute the improper integral
$$
\int_{-\infty}^{\infty} \frac{x+1}{x^4+2 x^2+5} d x
$$
\end{problem}

\begin{soln}
    Let $\pm\alpha$ and $\pm\beta$ be the roots of the denominator,
    which factors as
    \begin{equation*}
	(x^2 + 1)^2 + 4 = (x^2 + 2i + 1)(x^2 - 2i + 1),
    \end{equation*}
    and suppose also $\alpha^2 = -2i-1$ and $\beta^2 = 2i-1$
    and that $\alpha$ and $\beta$ have positive imaginary parts
    (relabel and/or $\alpha$ and $\beta$ otherwise).
\end{soln}



\begin{problem}
[2018 J-II-4]
Let $\left\{f_n: n=1,2,3, \ldots\right\}$ and $f$ denote holomorphic functions defined on an open domain $D$ in the complex plane. If the sequence $\left\{f_n\right\}$ converges uniformly on compact subsets of $D$ to $f,$ then show that the sequence of derivatives $\left\{f_n^{\prime}\right\}$ converges uniformly on compact subsets of $D$ to the derivative $f^{\prime}.$
\end{problem}

\begin{soln}
    This problem is the same as 2026 J-II-4.
\end{soln}



\begin{problem}
[2018 A-I-3]
Show that if $c>1,$ then the function
$$
f(z)=z e^{c-z}-1
$$
has precisely one root in $\Delta=\{|z|<1\},$ and this root is real and positive.
\end{problem}

\begin{soln}
    
\end{soln}



\begin{problem}
[2018 A-II-4]
Fix an infinite sequence $\left\{z_n\right\}$ of complex numbers in the unit disk $D$ and suppose that for every other sequence $\left\{w_n\right\}$ in $D,$ there is a bounded holomorphic function $f$ on $D$ so that $f\left(z_n\right)=w_n$ for all $n.$ Prove that there is a constant $C<\infty,$ so that we can always choose such an $f$ that also satisfies
$$
\sup _{z \in D}|f(z)| \leq C \sup _n\left|w_n\right| .
$$
\end{problem}

\begin{soln}

\end{soln}



\begin{problem}
[2017 J-I-5]
Is the ring of holomorphic functions on $\mathbb{C}$ a unique factorization domain? Justify your answer.
\end{problem}

\begin{soln}
    The answer is \framebox{no}.

    Let $f$ be an entire function which is irreducible
    in the ring $\mathcal O(\CC)$ of entire functions.
    If $f$ is nowhere vanishing then it is a unit,
    and if $f$ has at least two zeros (counting with multiplicity),
    say $z_1$ and $z_2$, then $f(z) = (z - z_1)(z - z_2)g(z)$
    for some entire function $g$,
    and $z-z_1$ and $(z-z_2)g(z)$ are non-units,
    so $f$ is reducible.
    Conversely if $f$ has exactly one zero $z_1$,
    then if $f = gh$ for non-units $g$ and $h$
    then $g$ and $h$ both have a zero,
    hence both are zero at $z_1$.
    This implies that $gh$ has a zero at $z_1$
    of order at least $2$, which contradicts the fact that
    $f$ has exactly one zero; thus $f$ is irreducible.
    That is, the irreducible elements of $\mathcal O(\CC)$
    are the entire functions with exactly one zero.

    Now observe that $\sin z\in\mathcal O(\CC)$
    has infinitely many zeroes
    while any finite product of irreducible elements
    of $\mathcal O(\CC)$ has finitely many zeroes,
    so $\sin z$ does not have a factorization into irreducibles,
    hence $\mathcal O(\CC)$ is not a UFD.
\end{soln}



\begin{problem}
[2017 J-II-5]
Let $\mathbb{D} \subset \mathbb{C}$ be the open unit disk and $\mathbb{H}:=\{z \in \mathbb{C}: \text{Im}(z)>0\}$ the upper half plane. Let $\Gamma$ be the set of holomorphic maps $f: \mathbb{D} \longrightarrow \mathbb{H}$ satisfying $f(0)=i.$ Compute
$$
\sup _{f \in \Gamma} \text{Im}(f(3 / 4))
$$
where $\text{Im}(z)$ is the imaginary part of $z$ for all $z \in \mathbb{C}.$
\end{problem}

\begin{soln}

\end{soln}



\begin{problem}
[2017 A-I-4]
Use the Residue Theorem to evaluate the integral $\int_0^{\infty} \frac{x^{1 / 2} d x}{x^3+1}.$
\end{problem}

\begin{soln}

\end{soln}



\begin{problem}
[2017 A-I-6]
Let $D \subset \mathbb{C}$ denote the open unit disk, and let $f: D \rightarrow \mathbb{R}$ be a harmonic function which belongs to $L^2(D).$ Prove that there is a positive constant $K$ such that
$$
|f(z)|<\frac{K}{1-|z|}
$$
for all $z \in D.$
\end{problem}

\begin{soln}

\end{soln}



\begin{problem}
[2016 J-I-4]
Let $\mathbb{D}=\{z:|z|<1\}$ be the unit disk in $\mathbb{C},$ and $f: \mathbb{D} \rightarrow \mathbb{D}$ be a holomorphic map such that $f^{\prime}(0)=0.$ Show that
$$
|f(z)-f(-z)| \leq 2|z|^3
$$
\end{problem}

\begin{soln}
    Define $g\colon\mathbb D\to\mathbb D$ by
    \begin{equation*}
        g(z) = \begin{cases}
	    (f(z) - f(-z))/2 & \text{if }z\ne 0\\
	    0 & \text{if }z = 0,
        \end{cases}
    \end{equation*}
\end{soln}



\begin{problem}
[2016 A-I-4]
Let $p: \mathbb{C} \longrightarrow \mathbb{C}$ be a polynomial with simple roots, none of which lies on the unit circle. Show that the number of roots inside the unit disk is equal to the degree of the map:
$$
\rho: S^1 \longrightarrow S^1, \quad \rho\left(e^{i \theta}\right)=\frac{p\left(e^{i \theta}\right)}{\left|p\left(e^{i \theta}\right)\right|}
$$
\end{problem}

\begin{soln}

\end{soln}



\begin{problem}
[2016 A-I-6]
Let $\Omega$ be a connected open subset of $\mathbb{C},$ and let $F$ be a meromorphic function on $\Omega.$ Show that there is a meromorphic function $h$ on $\Omega$ such that $h^{\prime}(z)=F(z)$ if and only if all the residues of $F$ vanish.
\end{problem}

\begin{soln}

\end{soln}



\begin{problem}
[2016 A-II-5]
Find all injective holomorphic maps $F: \mathbb{C}-\{0\} \longrightarrow \mathbb{C}-\{0\}.$
\end{problem}

\begin{soln}
    The solution set is
    $\boxed{F(z) = cz}$ and \framebox{$F(z) = \frac cz$}
    for some constant $c$.

    If $F$ has an essential singularity at $0$,
    then $\{f(z) : \abs z < \eps\}$ is dense in $\CC$
    for every $\eps > 0$ by the Casorati-Weierstrass theorem
    so [why?] $F$ cannot be injective.
    Thus $0$ is either removable or a pole;
    similarly, by considering $F(\frac 1z)$,
    which is also an injective holomorphic map
    $\CC - \{0\}\to\CC-\{0\}$ because $z\mapsto\frac 1z$
    is a bijection on $\CC - \{0\}$,
    we can conclude that $\infty$ is either removable or a pole.

    Thus, $F$ extends to a nonconstant map on the Riemann sphere.
    Such an extension must be surjective
    because the image of $F$ is open by the open mapping theorem
    and closed because it is (the image of) a compact set,
    which means that $F(\{0,\infty\}) = \{0,\infty\}$.
    By replacing $F$ with $\frac 1F$ if necessary
    we can assume that $F(\infty) = \infty$,
    so that $F$ is entire with a pole at $\infty$,
    hence a polynomial. By the fundamental theorem of algebra
    such a polynomial must have degree $1$
    (else there is a point in $\CC - \{0\}$
    whose preimage has $\deg F > 1$ points),
    so if $F(\infty) = \infty$ and $F(0) = 0$ then $F(z) = cz$
    for some constant $c$.
    The other case of $F(\infty) = 0$ and $F(0) = \infty$
    was reduced to the prior one by reciprocration,
    so the possible $F$ in this case are $F(z) = \frac cz$
    for some constant $c$.
\end{soln}



\begin{problem}
[2015 J-I-1]
Suppose $f$ and $g$ are two entire holomorphic functions with the property that $|f(z)| \leq|g(z)|$ for every $z \in \mathbb{C}.$ Prove that there is a constant $\lambda \in \mathbb{C}$ such that $f=\lambda g.$
\end{problem}

\begin{soln}
    If $z_0$ is a zero of $g(z)$ with order $d$,
    then $g(z) = (z-z_0)^d g_0(z)$ for some entire function $g_0$
    which is not zero at $z_0$.
    The given inequality implies $z_0$ is a zero of $f(z)$ as well,
    so write $f(z) = (z-z_0)^{d'}f_0(z)$
    for some entire function $f_0$ which is not zero at $z_0$.
    If $d' < d$, then
    \begin{equation*}
	\abs{f_0(z)} \le \smabs{(z-z_0)^{d - d'} g_0(z)},
    \end{equation*}
    which implies $f_0(z)$ is zero at $z_0$;
    this is a contradiction, so $d'\ge d$.
    That is, every zero of $g$ is a zero of $f$
    with higher order in $f$ than in $g$,
    so the singularities of $f/g$ at the zeroes of $g$
    are all removable.
    In other words $f/g$ extends to an entire function
    which by the given information satisfies $\abs{f(z)/g(z)}\le 1$
    for all $z$, whence Liouville's theorem implies
    $f/g = \lambda$ for some constant $\lambda\in\CC$.
\end{soln}



\begin{problem}
[2015 A-I-3]
Let $f$ be a holomorphic function defined on a neighborhood of the closed unit disk $\overline{\mathbb{D}}:=\{z \in \mathbb{C}:|z| \leq 1\},$ and assume that $|f(z)|<1$ for all $|z|=1.$ Determine the number of fixed points of $f$ in $\overline{\mathbb{D}}.$
\end{problem}

\begin{soln}

\end{soln}



\begin{problem}
[2015 A-I-5]
Prove that if $\left\{u_n\right\}$ is a sequence of harmonic functions on the unit disk with $\left|u_n(x)\right| \leq 1$ for all $x$ and all $n,$ then there is a subsequence which converges at every point of unit disk.
\end{problem}

\begin{soln}

\end{soln}



\begin{problem}
[2014 J-I-4]
Use the method of residues to compute the improper integral $\int_0^{\infty} \frac{x^2}{x^6+x^4+x^2+1} d x$
\end{problem}

\begin{soln}

\end{soln}



\begin{problem}
[2014 J-II-1]
Set $f(z)=\sum_{n=1}^{\infty} f_n(z),$ where $f_n(z)=\frac{z^n}{10^{n^2}}.$ Show that $f(z)$ is an entire function which has exactly ${n}$ distinct zeros inside of each circle $C_n=\left\{z:|z|=100^n\right\}, {n}=1,2,3,4, \ldots.$ . (Hint: compare the functions $\left|f_n(z)\right|$ and $\left|f(z)-f_n(z)\right|$ on $\left.C_n.\right)$
\end{problem}

\begin{soln}

\end{soln}



\begin{problem}
[2014 A-I-5]
Let $\mathbb{D} \subset \mathbb{C}$ be the unit disk, $f: \mathbb{D} \rightarrow \mathbb{D}$ a holomorphic function. Suppose that $f(c)=c$ for some $c \in \mathbb{D},$ and $\left|f^{\prime}(c)\right|<1.$ Prove that for any point $z_0 \in \mathbb{D},$ the sequence $\left\{z_n\right\}$ defined by iteration, $z_n=f\left(z_{n-1}\right),$ converges to $c.$
(Hint: first consider the case $c=0.$)
\end{problem}

\begin{soln}

\end{soln}



\begin{problem}
[2014 A-II-2]
Let $\Omega \subset \mathbb{C}$ be the right half-plane with the disk $D$ removed, where $D$ is the disk of radius $r=3$ centered at $w_0=5.$ What is the maximum number of disjoint open disks in $\Omega$ whose boundaries touch both boundary components of $\Omega$ ? Prove your answer.
(Hint: transform the domain.)
\end{problem}

\begin{soln}

\end{soln}



\begin{problem}
[2013 J-I-1]
What is the average of the angle $\alpha(z)$ subtended by the segment $[0,1]$ at a point $z$ in the disk $D=\{z \in \mathbb{C} : \abs{z-i} \leq 1\}$ ? Hint: express $\alpha(z)$ in terms of a holomorphic function on $D.$
\end{problem}

\begin{soln}

\end{soln}



\begin{problem}
[2013 J-II-2]
Show that the critical points of a complex polynomial $p(z)$ all lie in the closed convex hull of its roots. For the sake of simplicity, you may assume that the roots of $p(z)$ are all distinct.
\end{problem}

\begin{soln}

\end{soln}



\begin{problem}
[2013 A-I-3]
Suppose $f: D(0,1) \rightarrow D(0,1)$ is a holomorphic map from the open unit disc into the closed unit disc. Can the image of $f$ contain both the points 0 and 1 ?
\end{problem}

\begin{soln}

\end{soln}



\begin{problem}
[2013 A-II-1]
Write an explicit formula for the conformal map from a complex $z$-plane $\Delta$ with removed rays $-\infty < x < -1,1 < x < \infty$ to a strip $\Omega=\{0 < \operatorname{Im} w < 7\}.$ The map transforms the left ray to the bottom boundary component of $\Omega$ and the right ray to the top boundary component respectively.
\end{problem}

\begin{soln}

\end{soln}



\begin{problem}
[2012 J-I-3]
Let $p(z)$ be a complex polynomial whose critical values are all in the open, complex, unit disk $\Delta,$ i.e., $p^{\prime}(z)=0$ implies $p(z) \in \Delta.$ Prove that $p^{-1}(\Delta)$ is connected.
\end{problem}

\begin{soln}

\end{soln}



\begin{problem}
[2012 J-II-2]
Let $f$ be a continuous, complex-valued function on the closed disk $\bar{\Delta}=\{z \in$ $\mathbb{C}:|z| \leq 1\}$ which is holomorphic on the interior. Assume that $f$ restricted to the boundary $\partial \bar{\Delta}$ only takes values on the curve $\{x+i y \in \mathbb{C} \mid x, y \in \mathbb{R}, x y=1\}.$ Show that $f$ is constant on the whole disk.
\end{problem}

\begin{soln}

\end{soln}



\begin{problem}
[2012 A-I-4]
Find, with proof, the precise number of zeroes of the complex polynomial $p(z)=$ $z^9-2 z^6+z^2-8 z-2$ inside the annulus $1<|z|<2.$
\end{problem}

\begin{soln}

\end{soln}



\begin{problem}
[2012 A-II-6]
Denote by $(\mathcal{H},\langle\bullet, \bullet\rangle)$ the Hilbert space of holomorphic functions on the unit disk $D \subset \mathbb{C}$ which is the closure of the set of complex polynomial functions on the disk with respect to the Hermitian inner product
$$
\langle f(z), g(z)\rangle=\int_0^{2 \pi} f\left(e^{i \theta}\right) \overline{g\left(e^{i \theta}\right)} \frac{d \theta}{2 \pi} .
$$
Prove that the span of $\{1 /(z-n)\}_{n=2}^{\infty}$ is dense in $\mathcal{H}.$
\end{problem}

\begin{soln}

\end{soln}



\begin{problem}
[2011 J-I-1]
Prove that any three distinct points in the Euclidean plane $\mathbb{R}^2$ are the vertices of a topological triangle whose edges are arcs of Euclidean circles, such that the interior angle at each vertex is $90^{\circ}.$
\end{problem}

\begin{soln}

\end{soln}



\begin{problem}
[2011 J-I-4]
Let $u: \mathbb{C} \rightarrow \mathbb{R}$ be a $\mathscr{C}^2$-smooth, subharmonic function. Show that the function $\Psi: \mathbb{C} \rightarrow \mathbb{R}$ defined by
$$
\Psi(z):=\frac{1}{2 \pi} \int_0^{2 \pi} u\left(e^z e^{i \theta}\right) d \theta
$$
is subharmonic. Then show that the function $\mathbb{R} \ni x \mapsto \Psi(x+i 0)$ is convex and increasing.
\end{problem}

\begin{soln}

\end{soln}



\begin{problem}
[2011 A-I-3]
Let $\Omega \subset \mathbb{C}$ be an convex open set, and let $f: \Omega \rightarrow \mathbb{C}$ be a holomorphic function such that $\text{Re} f^{\prime}(z)>0$ for all $z \in \Omega.$ Prove that $f$ is injective.
\end{problem}

\begin{soln}

\end{soln}



\begin{problem}
[2011 A-II-3]
If $p(z)=z^n+a_1 z^{n-1}+\cdots+a_n$ is a monic polynomial, prove that
$$
\sup \{|p(z)| : |z| \leq 1\} \geq 1.
$$
\end{problem}

\begin{soln}
This problem is the same as 2025 A-I-3.
\end{soln}



\begin{problem}
[2010 J-I-1]
Let $f: \mathbb{C}-\{0\} \rightarrow \mathbb{C}$ be a holomorphic function on the punctured complex plane, and suppose that $f(2 z)=f(z)$ for all $z \neq 0.$ Prove that $f$ is constant.
\end{problem}

\begin{soln}
    The annulus $A = \{z\in\CC : 1 \le \abs z \le 2\}$
    is compact, so $f(A)$ is compact.
    Then, if $z\in\CC$ has magnitude $r > 0$,
    let $n$ be the integer such that $1\le 2^nr < 2$
    (precisely, $n = -\log_2r$), so $2^nz\in A$.
    We are given that $\abs{f(z)} = \abs{f(2^nz)}$,
    so $f(z)\in A$ for all $z\in\CC$.
    Thus the image of $f$ is a compact (hence closed) subset of $\CC$.
    If $f$ is nonconstant then its image
    is also open by the open mapping theorem
    (since the domain $\CC - \{0\}$ is open in $\CC$),
    which implies that the image is $\CC$
    because $\CC$ is connected --- but this is absurd
    since $\CC$ is not compact.
\end{soln}



\begin{problem}
[2010 J-II-6]
A fractional linear transformation maps the annulus $r<|z|<1($ where $r>0)$ onto the domain bounded by the two circles $\left|z-\frac{1}{4}\right|=\frac{1}{4}$ and $|z|=1.$ Find $r.$
\end{problem}

\begin{soln}

\end{soln}



\begin{problem}
[2010 A-I-4]
Prove that an entire function $f(z)$ on the complex plane is a polynomial if $\smabs{f(z^2)} \leq|f(z)|^2$ for every $z.$
\end{problem}

\begin{soln}

\end{soln}



\begin{problem}
[2010 A-II-3]
Let $f: \Delta \rightarrow \mathbb{C}$ be a holomorphic function on the open unit disk. Decompose $f$ into its real and imaginary parts, $f=u+i v.$ Among all holomorphic functions $f$ with $v(0)=0$ and with $|u|<1,$ what is the largest possible value of $v(1 / 2)$ ?
\end{problem}

\begin{soln}

\end{soln}



\begin{problem}
[2009 J-I-3]
Let $S=[-1,1]$ denote the closed line segment from $-1$ to 1 and suppose $f$ is the Riemann map from the unit disk to the domain
$$
\Omega=\{z: \operatorname{dist}(z, S)<1\} .
$$
What is the smoothness of $\partial \Omega$ and what is the radius of convergence of the power series of $f$ around the origin?
\end{problem}

\begin{soln}

\end{soln}



\begin{problem}
[2009 J-I-4]
Suppose that $u$ is a harmonic function on all of $\mathbb{R}^n$. Show that if $u \in L^1\left(\mathbb{R}^n\right)$, then $u \equiv 0$. Say $1< p<\infty$. Show that if $u \in L^p\left(\mathbb{R}^n\right)$, then $u \equiv 0$.
\end{problem}

\begin{soln}

\end{soln}



\begin{problem}
[2009 A-I-1]
Let $p(z)$ and $q(z)$ be two polynomials over the complex numbers $\mathbb{C}$, and suppose that their degrees satisfy $\text{deg}(p) \leq \text{deg}(q)-2$. Let $C_r \subset \mathbb{C}$ be the circle radius $r$ centered at 0 , and assume that each root $\zeta$ of $q(z)$ satisfies $|\zeta|< r$. Compute the contour integral
$$
\int_{C_r} \frac{p(z)}{q(z)} d z
$$
\end{problem}

\begin{soln}

\end{soln}



\begin{problem}
[2009 A-II-1]
Let $u: \mathbb{C} \rightarrow \mathbb{R}$ be a non-constant harmonic function on the complex plane. Show that the subset $\{z \in \mathbb{C} \mid u(z)=0\}$ is unbounded.
\end{problem}

\begin{soln}

\end{soln}



\begin{problem}
[2009 A-II-5]
Let $\Omega=\mathbb{C} \backslash\{a, b, c\}$ be the complement of three distinct points in the complex plane. Show that the set of one-to-one holomorphic maps $f: \Omega \rightarrow \Omega$ forms a finite group. Is the order of this group independent of the choice of $\{a, b, c\}$ ?
\end{problem}

\begin{soln}

\end{soln}



\begin{problem}
[2008 J-I-6]
Let $\Omega$ denote the open subset $\{x+i y \mid-1< y<1\}$ of the complex plane, and let $f: \Omega \longrightarrow \Omega$ denote a holomorphic function which satisfies $f(-1)=-1$ and $f(1)=1$. Show that $f$ is the identity map.
\end{problem}

\begin{soln}

\end{soln}



\begin{problem}
[2008 J-II-3]
Compute the real integral $\int_{-\infty}^{\infty} \frac{1}{1+x^4} d x$.
\end{problem}

\begin{soln}

\end{soln}



\begin{problem}
[2008 A-I-1]
Find all entire functions $f$ such that $\left|f^{\prime}\right| \leq|f|$ everywhere.
\end{problem}

\begin{soln}

\end{soln}



\begin{problem}
[2008 A-II-2]
Let $\Omega$ be the right half-plane $(\{x+i y: x>0\})$ with the disk of radius $r<1$ centered at 1 removed. Let $N(r)$ be the maximum number of disjoint open disks in $\Omega$ whose boundaries touch both boundary components of $\Omega$. Show
$$
N(r)=\left\lfloor\frac{\pi}{\arcsin \left(\sqrt{\frac{1-r}{1+r}}\right)}\right\rfloor .
$$
(Here $\lfloor\cdot\rfloor$ denotes the greatest integer function.)
\end{problem}

\begin{soln}

\end{soln}



\begin{problem}
[2007 J-I-2]
Find an explicit holomorphic bijection $\zeta=f(z)$ between the crescent
$$
\{z \in \mathbb{C}: |z|\leq 1,\ |z+1|\geq \sqrt{2}\}
$$
and the disk
$$
\{\zeta \in \mathbb{C}: |\zeta|\leq 1\}.
$$
\end{problem}

\begin{soln}

\end{soln}



\begin{problem}
[2007 J-II-1]
Let $f(z)$ be a continuous complex-valued function on the closed unit disk $D=\{z \in \mathbb{C}|| z \mid \leq 1\}$. If $f$ is holomorphic in the interior $|z|<1$ of $D$, and real-valued on the boundary circle $|z|=1$, show that $f$ must be constant.
\end{problem}

\begin{soln}

\end{soln}



\begin{problem}
[2007 A-I-6]
Let us consider a rational function $f(z)=z \frac{z-a}{1-\bar{a} z}$ of complex variable, where $|a|<1$. Show that $f$ restricts to a degree two covering map from the unit circle $\mathbb{T}$ to itself. Show that the Lebesgue measure $\mu$ on $\mathbb{T}$ is invariant under $f$, i.e.,
$$
\int \phi \circ f\, d \mu=\int \phi\, d \mu
$$
for any continuous function $\phi$ on $\mathbb{T}$.
(Hint: make use of the Dirichlet problem in the unit disk.)
\end{problem}

\begin{soln}

\end{soln}



\begin{problem}
[2007 A-II-1]
Let $\mathbf{D}=\{z:|z|<1\}$ be the unit disk, and let $f: \mathbf{D} \rightarrow \mathbf{D}$ be a holomorphic function such that $f(0)=\alpha \neq 0$. Prove that $f$ does not have zeros in the disk $\{z:|z|<|\alpha|\}$.
\end{problem}

\begin{soln}

\end{soln}



\begin{problem}
[2006 J-I-2]
Let
$$
f(z)=\sum_{n=0}^{\infty} z^{2^n}, \quad|z|<1
$$
Prove that for every $\zeta$ satisfying $\zeta^{2^l}=1$ for some $l \in \mathbb{N}$,
$$
\lim _{r \rightarrow 1^{-}} f(r \zeta)=\infty .
$$
\end{problem}

\begin{soln}

\end{soln}



\begin{problem}
[2006 J-I-4]
Let $p>2$ be a prime number. Find the absolute value of
$$
S=\sum_{n=0}^{p-1} \exp \left\{\frac{2 \pi i n^2}{p}\right\} .
$$
\end{problem}

\begin{soln}

\end{soln}



\begin{problem}
[2006 J-II-2]
Show that there is no meromorphic function satisfying
$$
f(z+1)=f(z), \quad f(z+i)=f(z)
$$
for all $z \in \mathbb{C}$, whose only poles are simple poles at the points $m+n i$, $m, n \in \mathbb{Z}$.
\end{problem}

\begin{soln}

\end{soln}



\begin{problem}
[2006 A-I-4]
Let $G=\{z \in \mathbb{C} \mid \text{Re} z>0$ and $\text{Im} z>0\}$, and let $f: G \rightarrow G$ be a non-trivial conformal map such that $f \circ f=i d$. Prove that $f$ has precisely one fixed point inside $G$.
\end{problem}

\begin{soln}

\end{soln}



\begin{problem}
[2006 A-II-1]
Find
$$
\lim _{n \rightarrow \infty} \int_0^{\infty}\left(\frac{x}{\sqrt{n}}\right)^{n^2} e^{-n x} d x
$$
\end{problem}

\begin{soln}

\end{soln}



\begin{problem}
[2006 A-II-4]
Show that $z^7-4 z^3+z-1$ has 3 zeros inside the unit circle (counted with multiplicity).
\end{problem}

\begin{soln}

\end{soln}



\begin{problem}
[2005 J-I-2]
Suppose $T$ is an equilateral triangle of side length 1 and let $f_{-}: T \rightarrow D=\{z:|z|<1\}$ be a conformal map which sends the center $c$ of the triangle to the origin in the disk (this is a 1-to-1, onto holomorphic map between the interiors which you may assume extends continuously to the boundary). What is the radius of convergence of the power series of $f$ around the point $c$ ?
\end{problem}

\begin{soln}

\end{soln}



\begin{problem}
[2005 A-II-1]
Let $\Omega=\mathbb{C} \backslash[-1,1]$. Give an example of a non-constant, bounded holomorphic function on $\Omega$. Show that such a function cannot extend continuously to all of $\mathbb{C}$.
\end{problem}

\begin{soln}

\end{soln}



\begin{problem}
[2004 J-I-3]
Suppose $\Omega=\{z \in \mathbb{C}:|z-1|>1$ and $|z+1|>1\}$. Let $U$ be a harmonic function on $\Omega$ with boundary values $U=0$ on $\{|z-1|=1\}$ and $U=1$ on $\{|z+1|=1, z \neq 0\}$. Find the value of $U(4)$.
\end{problem}

\begin{soln}

\end{soln}



\begin{problem}
[2004 J-I-6]
Let $g$ denote a holomorphic function on subset of the complex plane given by $|z|< r$, were $r$ is a fixed real number satisfying $r >1$. Suppose that $g(z) \leq 1$ holds for all $|z| \leq 1$. Show that for all $t \in \mathbb{C}$ with $|t|< 1$, the equation
$$
z=t g(z)
$$
has a unique solution $z=s(t)$ in the $\text{disc}|z|< 1$. Show that $t \rightarrow s(t)$ is a holomorphic function on the $\text{disc}|t|< 1$.
\end{problem}

\begin{soln}

\end{soln}



\begin{problem}
[2004 A-I-3]
Suppose $f$ is a holomorphic function in the open unit disc
$D=\{z \in \mathbb{C}|| z \mid<1\}$ such that $f(0)=0$ and $\left|f(z)+z f^{\prime}(z)\right|<1$ for all $z \in D$. Show that $|f(z)| \leq|z| / 2$ for all $z \in D$.
\end{problem}

\begin{soln}

\end{soln}



\begin{problem}
[2004 A-II-3]
Let $f, g \in \mathcal{O}(D) \cap C(\bar{D})$ where $D \subset \mathbb{C}$ is the open unit disc $\{z \in \mathbb{C} \mid$ $|z|<1\}$. Suppose there is a circle homeomorphism $h: \partial D \longrightarrow \partial D$ such that $f(t)=(g \circ h)(t)$ for all $t \in \partial D$. Prove that $f(\bar{D})=g(\bar{D})$; i.e. the two functions have the same range on the closed disc.
\end{problem}

\begin{soln}

\end{soln}



\begin{problem}
[2003 J-I-1]
Let $f$ be a holomorphic function on the punctured plane $0<|z|<\infty$. Assume that there exists a positive constant $C$ and a real constant $N$ such that
$$
|f(z)| \leq C|z|^N \text { for } 0<|z|<\frac{1}{2} .
$$
Show that $z=0$ is either a pole or a removable singularity for $f$ and find a bound for ord $\text{ro}_0 f$, the order of the function $f$ at 0 .
\end{problem}

\begin{soln}

\end{soln}



\begin{problem}
[2003 J-I-2]
State any form of Picard's Theorem. Let $n \geq 2$ be an integer. Show that there are no nowhere vanishing and nonconstant entire functions $f$ and $g$ that satisfy
$$
f^n+g^n=1 .
$$ Assume now that $n>2$. What are all the solutions $f$ and $g$ in the ring of entire functions that satisfy
$$
f^n+g^n=1 ?
$$
Hint: transform the problem to the setting of meromorphic functions on the complex line.
\end{problem}

\begin{soln}

\end{soln}



\begin{problem}
[2003 J-II-1]
Suppose $f$ is a meromorphic function on the plane which is holomorphic on the unit disk and has one simple pole only on the unit circle at the point $z_0$. Let $\sum_{n=0}^{\infty} a_n z^n$ be the power series for $f$ at the origin.
Show that $\lim _{n \rightarrow \infty} a_n z_0^n$ exists.
\end{problem}

\begin{soln}

\end{soln}



\begin{problem}
[2003 A-I-1]
Prove that for any $r>0$ there exists $N>0$ such that for every $n>N$, the polynomials
$$
P_n(z)=1+z+\frac{z^2}{2 !}+\ldots+\frac{z^n}{n !}
$$
have no zeroes in the disk $\{z \in \mathbf{C}:|z|< r\}$.
\end{problem}

\begin{soln}

\end{soln}



\begin{problem}
[2003 A-I-2]
Let $f$ be a holomorphic function on the unit disk $\{z \in \mathbf{C}:|z|<1\}$ such that $|f(z)|<1$ for $|z|<1$
If $f\left(\frac{1}{2}\right)=\frac{1}{3}$, find a sharp upper bound for $\left|f^{\prime}\left(\frac{1}{2}\right)\right|$.
\end{problem}

\begin{soln}

\end{soln}



\begin{problem}
[2003 A-II-1]
Show that
$$
\lim _{M \rightarrow \infty} \int_{S_M} \frac{e^{iz}}{z} d z=0,
$$
where $S_M$ is the semicircle in the upper half plane with center at 0 and radius $M>0$. Evaluate
$$
\int_0^{\infty} \frac{\sin x}{x} d x .
$$
\end{problem}

\begin{soln}

\end{soln}



\begin{problem}
[2002 J-I-1]
Suppose $\Omega$ is a horizontal strip $\{z: \text{Im} z \in(a, b)\}$. Suppose $f$ is holomorphic in $\Omega$ and $f(z+1)=f(z)$ for all $z \in \Omega$. Prove that $f$ has an expansion
$$
f(z)=\sum_{-\infty}^{\infty} c_n e^{2 \pi i n z}
$$
which converges uniformly in $\{\text{Im} z \in(a+\epsilon, b-\epsilon)\}$ for all $\epsilon>0$.
\end{problem}

\begin{soln}

\end{soln}



\begin{problem}
[2002 J-I-4]
Let $\left\{a_n\right\}_{n=0}^{\infty}$ be a sequence of complex numbers tending toward 0 . By definition, the infinite product $\prod_{n=0}^{\infty}\left(1+a_n\right)$ converges iff $\sum_{n=0}^{\infty} \log \left(1+a_n\right)$ converges to a finite value. The product is said to converge absolutely if the corresponding sum does. Show that $\prod_{n=0}^{\infty}\left(1+a_n\right)$ is absolutely convergent iff $\sum\left|a_n\right|<\infty$. Show that for $\text{Re}(s)>1$, the infinite product $\prod_{p \text { prime }} \frac{1}{1-p^{-s}}$ is absolutely convergent. For $\text{Re}(s)>1$, let $\zeta(s)=\sum_{n=1}^{\infty} n^{-s}$. Show that the infinite product from (ii) equals $\zeta(s)$ whenever $\text{Re}(s)>1$. (Hint: expand $\left(1-p^{-s}\right)^{-1}$ in a geometric series for each $p$.)
\end{problem}

\begin{soln}

\end{soln}



\begin{problem}
[2002 J-II-2]
Show that there is no entire function $f(z)$ with the property that, for each point $w$ in the plane, $f^{-1}(w)$ consists of exactly two points.
\end{problem}

\begin{soln}

\end{soln}



\begin{problem}
[2002 A-I-5]
Let $D$ be the domain in the extended complex plane $C \cup\{\infty\}$ exterior to the circles $|z-1|=1$ and $|z+1|=1$. Find a Riemann map (one-to-one holomorphic map) of $D$ onto the strip $S=\{z \in C ; 0<\text{Im} z<2\}$.
\end{problem}

\begin{soln}

\end{soln}



\begin{problem}
[2002 A-II-2]
Let II be the upper half-plane, and $f: \mathbb{E} \rightarrow \mathbb{H}$ an antiholomorphic map such that $f \circ f=I \mathrm{~d}$ Show that $f$ can be written as $f(z)=\overline{g(z)}$ where
$$
g(z)=(a z+b) /(c z+d), \quad a, b, c, d \text { real, } a d-b c=-1
$$ Show that $f$ is conjugate to $f_0: z \mapsto-\bar{z}$, i.e.,
$$
f=\phi \circ f_0 \circ \phi^{-1}
$$
for some biholomorphic map $\phi: I I \rightarrow$ II
\end{problem}

\begin{soln}

\end{soln}



\begin{problem}
[2002 A-II-4]
Let $\Omega$ be a connected and simply-connected region in $\mathbb{C}$ Show that for any real-valued function $u$ which is harmonic in $\Omega$ (i.e. $u \in \mathcal{C}^2(\Omega), u_{x x}+u_{y y}=0$ in $\Omega$, with subscripts denoting partial derivatives), there exists a real-valued function $v$ such that $u+i v$ is bolomorphic in $\Omega$. Give an example which shows that the conclusion of (i) can fail if $\Omega$ is not simply connected.
\end{problem}

\begin{soln}

\end{soln}



\begin{problem}
[2001 J-I-6]
Let $f(z)$ be a holomorphic function in $|z| \leq 1$ and $u=\text{Re} f$. Prove the following formula: for $|z|<1$,
$$
f(z)=\frac{1}{2 \pi} \int_0^{2 \pi} u(\zeta) \frac{\zeta+z}{\zeta-z} d \theta+\mathrm{i} C, \quad \zeta=e^{\mathrm{i} \theta}
$$
where $C$ is some real constant.
\end{problem}

\begin{soln}

\end{soln}



\begin{problem}
[2001 J-II-5]
Let $U$ be a domain in a complex plane such that its universal cover is the unit disk $D$ (that is, there is a map $\pi: D \rightarrow U$ which is holomorphic and is a covering map in the usual topological sense). Show that if $f$ is an entire function (i.e., a holomorphic function on $\mathbb{C}$ ) such that $f(z) \in U$ for all $z$, then $f$ is constant.
\end{problem}

\begin{soln}

\end{soln}



\begin{problem}
[2001 J-II-6]
Calculate
$$
\int_0^{\infty} \frac{\log x}{x^2+a^2} d x
$$
where $a \in \mathbb{R}, a>0$.
\end{problem}

\begin{soln}

\end{soln}



\begin{problem}
[2001 A-I-2]
Evaluate, with proof, $\lim _{t \rightarrow 0} \frac{1}{t} \int_0^1 \sin \left(\frac{t}{\sqrt{x}}\right) d x$.
\end{problem}

\begin{soln}

\end{soln}



\begin{problem}
[2001 A-I-5]
Let $\Omega=\{z=x+i y:|z|<1, y>0\}$, Find an explicit 1-1, onto holomorphic map of $\Omega$ to the unit disk.
\end{problem}

\begin{soln}

\end{soln}



\begin{problem}
[2001 A-II-5]
Given a finite sequence of points $\left\{a_k\right\}$ in the open unit disk $\{|z|<1\}$, consider the complex function
$$
B(z)=\prod_{k=1}^n \frac{z-a_k}{1-\overline{a_k} z} .
$$
Show that $B$ maps the unit circle $\{|z|=1\}$ into itself and that if $B(0)=0$, and $g$ is a continuous function on the circle, then $\int_0^{2 \pi}\left|g\left(e^{i \theta}\right)\right| d \theta=\int_0^{2 \pi}\left|g \circ B\left(e^{i \theta}\right)\right| d \theta$, (i.e., the composition operator $\rightarrow g \circ B$ is an isometry for the $L^1$ norm).
\end{problem}

\begin{soln}

\end{soln}



\begin{problem}
[2000 J-I-4]
Let $f(z)=\sum_{n=1}^{\infty} z^{2^n}$ for $|z|<1$. Show that $|f(z)| \leq C+C \log \frac{1}{1-|z|}$ for some $C<\infty$ and every $|z|<1$. (Hint: choose $N$ so that $2^{-N} \simeq 1-|z|$ and cut the series at $N$.)
\end{problem}

\begin{soln}

\end{soln}



\begin{problem}
[2000 A-I-1]
Let $D$ be the domain in the extended complex plane $\mathbb{C} \cup\{\infty\}$ exterior to the circles $|z-1|=1$ and $|z+1|=1$. Find a Riemann map (one-to-one holomorphic map) of $D$ onto the strip $S=\{z \in \mathbb{C} ; 0<\text{Im} z<2\}$.
\end{problem}

\begin{soln}

\end{soln}



\begin{problem}
[2000 A-I-6]
Show that $\lim _{M \rightarrow \infty} \int_{S_M} \frac{e^{\mathrm{i} z}}{z} d z=0$, where $S_M$ is the semi-circle in the upper half plane with center at 0 and radius $M>0$. Evaluate $\int_0^{\infty} \frac{\sin x}{x} d x$.
\end{problem}

\begin{soln}

\end{soln}



\begin{problem}
[2000 A-II-1]
Let $f$ be a holomorphic function on $|z|<1$ such that $|f(z)|<1$ for all $|z|<1$ and $f\left(\frac{1}{2}\right)=\frac{1}{3}$. Find a sharp upper bound for $\left|f^{\prime}\left(\frac{1}{2}\right)\right|$.
\end{problem}

\begin{soln}

\end{soln}



\begin{problem}
[2000 A-II-5]
Let $\mathbb{D}$ be the unit disk in the complex plane $\mathbb{C}$ endowed with the Lebesgue measure. Consider the Hilbert space $H=L^2(\mathbb{D})$. Show that the functions $z^n$ are orthogonal in $H$. Let $f(z)$ be an analytic function in some neighborhood of the closed disk $\overline{\mathbb{D}}$. Then
$$
\int_D|f(z)|^2 d x d y=\pi \sum_{n=0}^{\infty} \frac{1}{n+1}\left|a_n\right|^2,
$$
where $a_n$ are the Taylor coefficients of $f$ at the origin.
\end{problem}

\begin{soln}

\end{soln}



\begin{problem}
[1999 J-I-3]
Suppose $f$ is an entire function so that $f^{-1}(K)$ is a compact set for every compact $K$ in the complex plane. Prove that $f$ is a polynomial.
\end{problem}

\begin{soln}

\end{soln}



\begin{problem}
[1999 J-II-6]
Suppose $D_0\subset \mathbb{R}^2$ is a disk centered at the origin and $\{D_k\}$, $k=1,\ldots,n$, are disjoint closed subdisks centered on the real axis. Let $u$ be the harmonic function in
$$
\Omega=D_0\backslash\bigcup_k D_k
$$
with boundary values $1$ on $\partial D_0$ and $0$ on the $\partial D_k$, $k=1,\ldots,n$. Prove that $u$ has exactly $n-1$ critical points, and all of them are on the real axis.
\end{problem}

\begin{soln}

\end{soln}



\begin{problem}
[1999 A-I-3]
Suppose $D=\{|z|<1\}$ is the unit disk and $\Omega$ is a simply connected domain. Suppose $f$ and $g$ are holomorphic maps of $D$ into $\Omega$ such that $f(0)=g(0)$ and that $f$ is also 1-1 and onto. Prove that $g(D(0,r))\subset f(D(0,r))$ for every $0< r<1$.
\end{problem}

\begin{soln}

\end{soln}



\begin{problem}
[1999 A-II-5]
Let $g$ be a holomorphic function on $|z|< R$, $R>1$, with $|g(z)|<1$ for all $|z|<1$. Show that for all $t\in\mathbb{C}$ with $|t|<1$, the equation
$$
z=tg(z)
$$
has a unique solution $z=s(t)$ in the disc $|z|<1$. Show that $t\mapsto s(t)$ is a holomorphic function on the disc $|t|<1$. (Hint: find an integral formula for $s$.)
\end{problem}

\begin{soln}

\end{soln}



\begin{problem}
[1998 J-I-3]
Let $f:S\rightarrow S$ be a conformal automorphism of the square $-3<\text{Re}z,\text{Im}z<3$ with $f(0)=2+2i$. Show that
$$
\frac{\sqrt2}{6}<|f'(0)|<\frac{10\sqrt2}{6}.
$$
\end{problem}

\begin{soln}

\end{soln}



\begin{problem}
[1998 J-I-4]
Let $f$ be a non-constant periodic entire function with $f(z+1)=f(z)$. Show that $f$ has a fixed point.
\end{problem}

\begin{soln}

\end{soln}



\begin{problem}
[1998 J-II-4]
Let $n$ be an integer, $n>1$. Let $h:D\rightarrow\mathbb{C}$ be a harmonic function on the unit disk which has boundary values $1$ on the segment $\{e^{i\theta}\mid \theta\in(0,2\pi/n)\}$ and has boundary values $0$ on the exterior of this arc. Show that $h(0)=1/n$.
\end{problem}

\begin{soln}

\end{soln}



\begin{problem}
[1998 A-I-2]
In $\mathbb{C}$, let $D=\{z:|z|<1\}$, $H_+=\{z:\text{Im}z>0\}$, $H_-=\{z:\text{Im}z<0\}$, $S=\{z:z\in\mathbb{R}\text{ and }-1< z<1\}$. Suppose that $g$ is continuous on the closure of $D$ and is holomorphic in $D$. For $z\in\mathbb{C}\backslash S$, let
$$
F(z)=\frac{1}{2\pi i}\int_{-1}^{1}(z-t)^{-1}g(t)\,dt.
$$ Show that $F$ is holomorphic in $\mathbb{C}\backslash S$. Show that $F$, restricted to $H_+$, has a holomorphic extension to $H_+\cup D$; call that extension $F_+$. Similarly, $F$, restricted to $H_-$, has a holomorphic extension to $H_-\cup D$; call that extension $F_-$. On $D$, consider the function $G=F_+-F_-$. Express $G$ in terms of $g$ without using any integrals.
\end{problem}

\begin{soln}

\end{soln}



\begin{problem}
[1998 A-I-6]
Let $p(z)=\sum_{k=0}^{n}a_kz^k$ be a polynomial over $\mathbb{C}$, with $a_0=a_n=1$. Let $\omega$ be a primitive $n$th root of unity, so that the $n$th roots of unity are $\{1,\omega,\omega^2,\ldots,\omega^{n-1}\}$. Prove that $\sum_{j=0}^{n-1}p(\omega^j)=2n$. Suppose that whenever $0< j< n$, $|p(\omega^j)|<2$. Show that $p(z)=z^n+1$ for all $z$.
\end{problem}

\begin{soln}

\end{soln}



\begin{problem}
[1998 A-II-2]
In $\mathbb{C}$, let $L=\{z:z\in\mathbb{R}\text{ and }z<0\}$.
Show that there exists a holomorphic function $h$ on $\mathbb{C}\backslash L$ such that for all $z=re^{i\theta}\in\mathbb{C}\backslash L$ (with $r>0$, $-\pi<\theta<\pi$), we have $|h(z)|=r^\theta$.
(Hint: in polar coordinates, the Laplacian on $\mathbb{R}^2$ is $\Delta=r^{-1}\partial_r(r\partial_r)+r^{-2}\partial_\theta^2$.)
\end{problem}

\begin{soln}

\end{soln}



\begin{problem}
[1997 J-I-3]
Let $n$ be a non-negative integer. What is the most general complex-analytic function $f:\mathbb{C}-\{0\}\rightarrow\mathbb{C}$ such that $f(2z)=2^n f(z)$ for all $z\in\mathbb{C}-\{0\}$? Explain.
\end{problem}

\begin{soln}

\end{soln}



\begin{problem}
[1997 J-II-1]
Let $f(z)$ be a continuous complex-valued function on the closed unit disc $|z|\leq 1$ which is holomorphic in the open disc $|z|<1$. Suppose that $f(e^{i\theta})=0$ for $\theta\in[0,\pi/2]$. Show that $f\equiv 0$.
\end{problem}

\begin{soln}

\end{soln}



\begin{problem}
[1997 A-I-3]
Let $\Sigma\subset\mathbb{C}$ be the unit square $[0,1]\times[0,1]$, and let $P\subset\mathbb{C}$ be the rectangle $[0,1]\times[0,2]$. Is there a homeomorphism $f:\Sigma\rightarrow P$ which is complex-analytic in the interior of $\Sigma$ and takes the corners of $\Sigma$ to the corners of $P$? Explain.
\end{problem}

\begin{soln}

\end{soln}



\begin{problem}
[1997 A-II-4]
Let $U$ be a connected open subset of the complex plane $\mathbb{C}$, and suppose that $f:\mathbb{C}\rightarrow U$ is both an entire analytic function and a covering map. Show that the fundamental group of $U$ is Abelian.
\end{problem}

\begin{soln}

\end{soln}



\begin{problem}
[1997 A-II-5]
Let $z=x+iy$ be the usual complex coordinate on $\mathbb{C}$, and let $\Sigma\subset\mathbb{C}$ be the unit square $0\leq x\leq 1$, $0\leq y\leq 1$. Suppose that $f:\Sigma\rightarrow\mathbb{C}$ is a smooth function which satisfies $f|_{\partial\Sigma}\equiv 0$. Prove that
$$
\int_\Sigma\left|\frac{\partial f}{\partial \bar z}\right|^2dx\,dy\geq \frac{\pi^2}{2}\int_\Sigma |f|^2dx\,dy.
$$
When does equality occur?
\end{problem}

\begin{soln}

\end{soln}



\end{document}
